%------------------------------------------------------------------------------%
\chapter{Funções Exponenciais}

%------------------------------------------------------------------------------%
\section{Potências}

\begin{defin}\label{def:exp}
  Seja $a$ um \textbf{número real positivo} chamado \textbf{base}. Definimos
  para todo $n\in\N$,
  \[
    a^n:=a.a.a.\cdots.a \quad~~~(n\text{ vezes}).
  \]
\end{defin}

Em particular, $a^1=a$. Segue da Definição \ref{def:exp} que para todo
$m, n\in\N$ as seguintes propriedades são válidas:
\begin{description}
  \item[P.1] $a^m.a^n=a^{m+n}.$
  \item[P.2] $(a^m)^n=a^{m.n}.$
  \item[P.3] Se $b>0$, então $(ab)^m=a^m.b^m.$
\end{description}

Fixada uma base $a$, queremos estender a Definição \ref{def:exp} para o
conjunto dos números reais, ou seja,  definir $a^x$ para todo $x\in\R$, de tal
forma que as propriedades acima continuem válidas.
Faremos essa extensão gradualmente:
\footnote{Observe que fixamos uma base $a$,
  ou seja, $a$ por definição é maior que zero. Posteriormente falaremos do porquê
  dessa escolha.
}
\begin{enumerate}
  \item Como definir $a^0$?

    Fazendo $m=0$ e $n=1$ em \textbf{P.1} temos que
    \[
      a^0.a^1=a^{0+1}=a^1
    \]
    Como $a\neq 0$, dividimos ambos os lados da equação por $a=a^1$ e
    concluímos que $a^0=1$ . Portanto, definimos $a^0:=1$. 
    \footnote{
      Observe que $0^0$ é uma expressão não definida, ou seja, $0^0$ não é um
      número real. 
    }

  \item Como definir $a^{-m}$ para $m$ inteiro positivo?

    Novamente aplicando a propriedade \textbf{P.1} e utilizando que $a^0=1$,
    temos que \[a^{-m}\cdot a^{m}=a^{-m+m}=a^0=1.\] Portanto, definimos
    \begin{equation}
      \label{def: nun neg}
      a^{-m}:=\frac{1}{a^m}.
    \end{equation}

    \begin{exe}
    \begin{enumerate}
      \item $3^{-2}=\frac{1}{3^2}=\frac{1}{9}.$
      \item $\left(\frac{2}{3}\right)^{-3}
        =\frac{1}{(\frac{2}{3})^3}
        =\frac{1}{\frac{2^3}{3^3}}
        =\frac{3^3}{2^2}
        =\left(\frac{3}{2}\right)^{3}
        =\frac{27}{8}$.
    \end{enumerate}
    \end{exe}

    Observem que se $a$ e $b$ são números positivos, então
    \[
      \left(\frac{a}{b}\right)^{-n}=\left(\frac{b}{a}\right)^{n}.
    \]

  \item Como definir $a^x$ para $x\in\mathbb{Q}$?

    Inicialmente faremos um exemplo: Considerando $x=n/2$, segue de
    \textbf{P.2} que
    \[ 
      a^{n/2}.a^{n/2}=(a^{n/2})^2=a^{2\cdot n/2}=a^n.
    \]
    Como o produto de $a^{n/2}$ por ele mesmo resultou em $a^n$, concluímos que
    $ a^{n/2}=\sqrt{a^n}.$

    De forma análoga, a propriedade \textbf{P.2} pode ser utilizada para
    definir $a^{n/m}$ para todo $n/m\in\mathbb{Q}$, a partir da função
    $f(x)=\sqrt[m]{x}$. Como $a^n$ já está definida para todo $n\in\mathbb{N}$
    e ao multiplicarmos $a^{n/m}$ por si mesmo $m$ vezes, obtemos
    \[ 
      a^{n/m}.\cdots.a^{n/m}=(a^{n/m})^m=a^{m \cdot n/m}=a^n,
    \]
    temos que $a^{n/m}:=\sqrt[m]{a^n}$.

    \begin{exe}
    \begin{enumerate}
      \item $3^{\frac{5}{2}}=\sqrt{3^5}=\sqrt{3^4\cdot 3}=9\sqrt{3}$.
      \item $\left(\frac{2}{3}\right)^{-\frac{2}{3}}
        =\left(\frac{3}{2}\right)^{\frac{2}{3}}
        =\sqrt[3]{\frac{9}{4}}$.
    \end{enumerate}
    \end{exe}

  \item Como definir $a^x$ para $x\in\mathbb{I}$?

    Fixado $x$ irracional, existe uma sequência de números racionais $x_n$ tal
    que os valores dessa sequência ficam cada vez mais próximos do número
    irracional $x$. Por exemplo, se $x=\pi=3,141592653589...$, podemos
    considerar a sequência dada por $x_1=3$; $x_2=3,1$; $x_3=3,14$;
    $x_4=3,141$; $x_5=3,1415$; $x_6=3,14159$; $x_7=3,141592;\cdots$. Nesse
    caso, dizemos que o limite quando $n$ tende a infinito da sequência $x_n$
    converge para $x$ (estudaremos com mais detalhe a ideia de limite no
    próximo módulo, nesse momento, apenas note que quando $n$ é muito grande o
    número $x_n$ fica muito próximo do número irracional $x$) e escrevemos:
    \[ 
      \lim_{n\rightarrow\infty}x_n=x. 
    \]

    Para cada $x_n$, $a^{x_n}$ já foi definido. A sequência $y_n=a^{x_n}$
    converge para um valor (ou seja, à medida que $n$ cresce, os valores
    $a^{x_n}$ ficam muito próximos de um número real que chamaremos de $a^x$).
    Definimos então
    \[
      a^{x}=\lim_{n\rightarrow\infty}a^{x_n}. 
    \]
\end{enumerate}

%------------------------------------------------------------------------------%
\section{Potências de números negativos}

Vimos acima que é possível definir $a^x$ para todo $a>0$ e $x\in\mathbb{R}$.
Por que não definimos para $a\leq 0?$ Faz sentido falar em $a^x$ para $a\leq
0$? A resposta é positiva se $x$ é um número inteiro . Se $x$  não é inteiro,
essa definição não pode ser feita sem gerar inconsistências, ou seja, não é
possível definir. Por exemplo, $(-4)^{\frac{1}{2}}$ não está definido, pois a
definição natural para esse número seria $\sqrt{-4}$ que não é um número real.
Aparentemente, poderíamos definir $(-8)^{\frac{1}{3}}$ como $\sqrt[3]{-8}=-2$.
Mas há um problema que surge da equivalência de frações: Sabemos que
$\frac{1}{3}=\frac{2}{6}$, logo teríamos que
$(-8)^{\frac{1}{3}}=(-8)^{\frac{2}{6}}$ mas
$(-8)^{\frac{2}{6}}=\sqrt[6]{(-8)^2}=\sqrt[6]{2^6}=2$.

Ou seja, uma potência de base negativa e expoente racional não pode ser
definida. Consequentemente, a definição de $a^x$ para $x$ irracional também não
está definida para $a<0$. Portanto $(-2)^{\frac{1}{2}},(-3)^{\sqrt{3}}$ não são
números reais.

%------------------------------------------------------------------------------%
\section{Funções Exponenciais}

Antes de definirmos as funções exponenciais, apresentaremos um exemplo para
termos noção como essas funções aparecem no dia a dia.

\begin{exe}
  \label{ex 1}
  Heloísa contraiu um empréstimo de $R\$1.000,00$ em um banco que cobra $6\%$
  de juros ao mês. Se Heloísa não saldar sua dívida, quanto ela deverá ao banco
  após quatro meses? E após um número $x$ arbitrário de meses?

  Observe que após um mês que o empréstimo foi contraído a dívida de Heloísa é
  
  \begin{align*}
    1000+(1000\cdot \frac{6}{100})
    & = 1000+1000\cdot 0,06\\
    & = 1000\cdot(1+0.06)\\
    & = 1000\cdot 1,06.
  \end{align*}
  Após dois meses, a dívida é
  \[
    (1000\cdot 1,06)\cdot 1,06 = 1000 \cdot (1,06)^2
  \]
  Após três meses, a dívida é
  \[
    (1000\cdot (1,06)^2)\cdot 1,06 = 1000 \cdot (1,06)^3
  \]
  E finalmente, após quatro meses, a dívida é
  \[
    (1000\cdot (1,06)^3)\cdot 1,06 = 1000 \cdot (1,06)^4\approx 1.262,48
  \]
  Generalizando a conta acima, após $x$ meses a dívida de Heloísa é
  \[
    1000 \cdot (1,06)^x.
  \]
\end{exe}

\begin{defin}
  \label{def:exp:1}
  A função exponencial com  base $a$ é definida por
  
  \begin{align}
  &f:\R\rightarrow\R \nonumber \\
  &x\rightarrow a^x
  \end{align}
  em que $a>0, a\neq 1$ é um número real fixo chamado base e $x$ é qualquer
  número real.
  \footnote{
    Para $a=1$, a função $f(x)=1^x$ é a função constante
    $f(x)=1$. Para $a\leq 0$, vimos na seção anterior que não faz sentido essa
    definição.
  }
\end{defin}

Segue que a função assim definida satisfaz as seguintes propriedades: Para todo
$x,y\in\R$, 
\footnote{
  As três primeiras propriedades estão listadas acima para
  os naturais, aqui estendemos para os reais.
}
\begin{description}
  \item[P.1] $ a^x.a^y=a^{x+y}.$
  \item[P.2] $(a^x)^y=a^{x.y}.$
  \item[P.3] Se $b>0$, então $(ab)^x=a^x.b^x.$
  \item[P.4] $\frac{a^x}{a^y}=a^{x-y}.$
\end{description}

A função exponencial pode ser classificada em crescente ou decrescente, a
depender do valor de $a$:
\begin{description}
  \item[1:] Se $a>1$ a função $f(x)=a^x$ é
    crescente\footnote{$x<y\Leftrightarrow a^x<a^y$} e o seu gráfico tem o
    formato das curvas a seguir:

  \begin{figure}[h]
    \centering
    \includegraphics[width=0.45\textwidth,trim={0 1cm 9cm 3cm},clip]{exp1}
    \caption{ Gráfico de $f(x)=2^x$ (traço contínuo) e $g(x)=4^x$ (tracejato).}
  \end{figure}
  
  \item[2:] Se $0<a<1$ a função $f(x)=a^x$ é
    descrescente\footnote{$x<y\Leftrightarrow a^x>a^y$} e o seu gráfico tem o
    formato das curvas a seguir:
  
  \begin{figure}[h]
    \centering
    \includegraphics[width=0.45\textwidth,trim={0 1cm 9cm 3cm},clip]{exp2}
    \caption{ Gráfico de $f(x)=2^{-x}=\big(\frac{1}{2}\big)^x$ (traço
      contínuo) e $g(x)=4^{-x}=\big(\frac{1}{4}\big)^x$ (tracejato).}
  \end{figure}
\end{description}
Note que, uma vez que $a^0=1$, para todo $a\neq 0$, então o gráfico da função
$a^x$ sempre intercepta o eixo $y$ no ponto (0,1).

%------------------------------------------------------------------------------%
\subsection{Procedimento para esboçar o gráfico da função exponencial}

Se quisermos esboçar o gráfico da função $f(x)= a^x$, escolhemos um conjunto de
valores (3 ou 5 pontos) estratégicos para $x$ e montamos uma tabelinha com
esses valores e sua imagem pela função. É comum escolhermos os valores
$x=-2,-1,0,1,2$. Vamos observar os seguintes exemplos.

\begin{exe}
  Nesse exemplo, vamos fazer o esboço do gráfico da função $f(x)=2^x$. Para
  isso, vamos escolher os valores $x=-2,-1,0,1,2$ e montar uma tabela com 3
  colunas: A primeira coluna esses valores de $x$, a segunda coluna com a
  imagem desses valores pela função $f(x)=2^x$ e a terceira com o par ordenado
  $(x,f(x))$. A tabela \ref{2el} mostra esses valores e a figura \ref{2eleva}
  mostra o gráfico de $f$.
\end{exe}

\begin{table}
    \centering
    \caption{Tabela de valores para o gráfico de 
      $f(x)=\left(\frac{1}{2}\right)^x$.}
    \label{2el}
    \begin{tabular}{|c|c|c|}
        \hline
        $x$ & $f(x)=2^x$ & $(x,f(x))$  \\
        \hline
        -2 & $f(-2)=2^{-2}=\frac{1}{4}$ & $(-2,\frac{1}{4})$ \\
        \hline
        -1 & $f(-1)=2^{-1}=\frac{1}{2}$ & $(-1,\frac{1}{2})$ \\
        \hline
        0 & $f(0)=2^0=1$ & $(0,1)$ \\
        \hline
        1 & $f(1)= 2^1=2$ & $(1,2)$  \\
        \hline
        2 & $f(2)=2^2=4$ & $(2,4)$  \\
        \hline
    \end{tabular}
\end{table}

\begin{figure}[h!]
    \centering
    \includegraphics{2elevado-1}
    \caption{Esboço do gráfico da função $y=2^x$.}
    \label{2eleva}
\end{figure}

\begin{exe}
  Nesse exemplo, vamos fazer o esboço do gráfico da função $f(x)=\left(
  \frac{1}{2}\right)^x$. Para isso, vamos escolher, novamente, os valores
  $x=-2,-1,0,1,2$ e construir uma tabela com 3 colunas: A primeira coluna esses
  valores de $x$, a segunda coluna com a imagem desses valores pela função
  $\left( \frac{1}{2}\right)^x$ e a terceira com o par ordenado $(x,f(x))$. A
  tabela \ref{2meioel} mostra esses valores e a figura \ref{meioeleva} mostra o
  gráfico de $f$.
\end{exe}

\begin{table}[h!]
    \centering
    \caption{Tabela de valores para o gráfico de 
      $f(x)=\left(\frac{1}{2}\right)^x$.}
    \label{2meioel}
    \begin{tabular}{|c|c|c|}
        \hline
        $x$ & $f(x)=\left(\frac{1}{2}\right)^x$ & $(x,f(x))$  \\
        \hline
        -2 & $f(-2)=\left(\frac{1}{2}\right)^{-2}=2$ & $(-2,2)$ \\
        \hline
        -1 & $f(-1)=\left(\frac{1}{2}\right)^{-1}=2$ & $(-1,1)$ \\
        \hline
        0 & $f(0)=\left(\frac{1}{2}\right)^0=1$ & $(0,1)$ \\
        \hline
        1 & $f(1)= \left(\frac{1}{2}\right)^1=\frac{1}{2}$ & $(1,\frac{1}{2})$  \\
        \hline
        2 & $f(2)=\left(\frac{1}{2}\right)^2=\frac{1}{4}$ & $(2,\frac{1}{4})$  \\
        \hline
    \end{tabular}
\end{table}

\begin{figure}[h!]
    \centering
    \includegraphics{meioelevado}
    \caption{Esboço do gráfico da função $y=\left(\frac{1}{2}\right)^x$.}
    \label{meioeleva}
\end{figure}

%------------------------------------------------------------------------------%
\section{Equações e Inequações Exponenciais}

%------------------------------------------------------------------------------%
\subsection{Equações exponenciais}

As equações em que a incógnita aparece nos expoentes são chamadas de
\textit{Equações exponenciais}. A seguir serão apresentados alguns exemplos de
equações exponenciais, juntamente com as resoluções. Para resolvê-las usaremos
o fato de que a função exponencial é injetora, isto é:
\[
    a^{x_1}=a^{x_2} \Rightarrow x_1=x_2.
\]

\begin{exe}
    Resolva as seguintes equações  exponenciais:
    \begin{enumerate}
      \item $\left( \frac{1}{3}\right)^x=81$ 

        A ideia para resolver uma equação exponencial desse tipo, tranformamos
        a equação numa igualdade de potências de mesma base:
        \[
          \left( \frac{1}{3}\right)^x=81 \Rightarrow 3^{-x}=3^4
        \]
        Como a função exponencial é injetora então basta igualar os expoentes,
        o que implica que $-x=4 \Rightarrow x=-4$. Logo, $S=\{ -4 \}$.

      \item $5^x-125=0$ \\

        Podemos reescrever a equação como $5^x=5^3$, ou seja, transformando em
        uma igualdade de potências de mesma base. Como a função exponencial é
        injetora então basta igualar os expoentes, o que implica que $x=3$.
        Logo, $S=\{ 3 \}$.

      \item $ \left( \frac{1}{2} \right)^x=\sqrt[3]{4}$ 

        Note que
        \[ 
          \left( \frac{1}{2} \right)^x
          = \sqrt[3]{4} \Rightarrow 2^{-x}
          = (2^2)^{\frac{1}{3}} \Rightarrow 2^{-x}
          = 2^{\frac{2}{3}} \Rightarrow -x
          = \frac{2}{3} \Rightarrow x=-\frac{2}{3}
        \]
        Logo, $S=\{ \frac{-2}{3} \}$

      \item $2^{x-3}+2^{x-1}+2^x=52$ 

        Nesse caso, não conseguimos transformar a equação em uma igualdade de
        mesma base. Desse modo, vamos reescrever a equação, usando as
        propriedades de potenciação:
        \[
          2^x. 2^{-3}+2^x .2^{-1}+ 2^x=52
        \]
        Colocando $2^x$ em evidência
        \[
          2^x \left( \frac{1}{8}+ \frac{1}{2}+1\right)
          =52 \Rightarrow 2^x. \frac{13}{8}
          =52 \Rightarrow 2^x
          = 32 
        \]
        Portanto, colocando $32$ na base $2$ temos $2^x=2^5 \Rightarrow x=5$.
        Logo, a solução da equação é $S=\{5 \}$.
    
      \item $2^{2x}-9.2^x+8=0$ 

        Do mesmo modo do item anterior, , não conseguimos transformar a equação
        em uma igualdade de mesma base. Vamos reescrever a equação da seguinte
        forma:
        \[
          (2^x)^2-9.2^x+8=0
        \]
        Agora, vamos fazer a seguinte mudança de variável $y=2^x$. Assim,
        \[
          y^2-9y+8=0
        \]
        Que é uma equação do segundo grau, cujas as raízes são $y=1$ e $y=8$.
        Agora, precisamos encontrar os valores de $x$ correspondentes.
        \begin{enumerate}
            \item Se $y=1 \Rightarrow 2^x=1 \Rightarrow 2^x=2^0 \Rightarrow x=0$.
            \item Se $y=8 \Rightarrow 2^x=2^3 \Rightarrow x=3$.
        \end{enumerate}
        Logo o conjunto solução da equação é $S=\{0,3 \}$.
    \end{enumerate}
\end{exe}

%------------------------------------------------------------------------------%
\subsection{Inequações Exponenciais}
    
Para resolvermos uma Inequação exponencial do tipo $a^{f(x)} \leq a^{g(x)}$ (ou
$a^{f(x)} \geq a^{g(x)}$), devemos lembrar que a função $y=a^x$ é crescente se
$a>1$ e decrescente se $0 < a< 1$. Portanto, antes de tudo, observamos a
situação das bases dos dois membros e, caso as bases sejam diferentes,
reduza-as as mesmas bases e depois monte uma inequação respeitando as seguintes
condições:
\begin{enumerate}
    \item Se $a>1$ conserve o sinal original da inequação. Isso se dá pelo fato
      da funçãio exponencial ser crescente se $a>1$.

    \item Se $0< a< 1$ inverta o sentido da desigualdade. Essa fato, se dá por
      conta da função exponencial ser decrescente quando $0< a< 1$.
\end{enumerate}

\begin{exe} 
  Resolva as seguintes Inequações Exponenciais:
  \begin{enumerate}
    \item $\left( \frac{1}{2} \right)^{2x+2} 
      < \left( \frac{1}{2} \right)^{x^2-1} $ 

      Note que nesse caso temos uma inequação exponencial e como as bases de ambos
      os membros são iguais( $a=\frac{1}{2}$) basta montarmos uma inequação com
      os expoentes, invertendo o sentido da desigualdade, visto que $0 < a <
      1$. Assim, precisamos resolver a seguinte inequação exponencial $2x+2>
      x^2-1 \Rightarrow x^2-2x-3< 0$, que é uma inequação do segundo grau, cuja
      solução é $S= \{ x \in \mathbb{R} | -1<x<3\}$.

  \item $2^{x+2} \geq 64$

      Inicialmente, observamos que $64=2^6$. Portanto, estamos interessado em
      resolver a inequação $2^{x+2} \geq 2^6$. Como já deixamos ambos os
      membros da desigualde na mesma base então montamos uma inequação com os
      expoentes mantendo o sentido da desigualdade original, uma vez que a base
      $a=2>1$.

      Assim, temos que resolver a seguinte inequação do primeiro grau $x+2 \geq 6
      \Rightarrow x> 4$. Logo, a solução da inequação pedida é $S= \{ x \in
      \mathbb{R}| x>4 \}$.
  \end{enumerate}
\end{exe}

\begin{exe}
  Determine o domínio D da função $f(x)= \sqrt{5^x-125}$.
  
  A fim de que exista $f(x)$, devemos ter $5^x-125 \geq 0$, uma vez que, não
  existe raíz quadrada de número negativo. Assim,
  \[
    5^x-125 \geq 0 
    \Rightarrow 5^x \geq 125 
    \Rightarrow 5^x \geq 5^3
    \Rightarrow x \geq 3
  \]
  Assim, o domínio da função é D=$\{ x \in \mathbb{R}| x \geq 3 \}.$
\end{exe}

%------------------------------------------------------------------------------%

%------------------------------------------------------------------------------%
\section{Exercícios}

\begin{exercises*}

%------------------------------------------------------------------------------%
\begin{exercise}
  Faça o esboço gráfico das funções abaixo
  \begin{mathlist}
    \item f(x) = 2^x-2
    \item f(x) = \left( \frac{1}{5} \right)^x
    \item f(x) =-4\left( \frac{1}{2} \right)^x
    \item f(x) = 3^{-x}
  \end{mathlist}

  \begin{answer}
    \begin{alphalist}
        \item \IncludeGraphicCentred{gabarito_aula4-1a}[0.9]
        \item \IncludeGraphicCentred{gabarito_aula4-1b}[0.9]
        \item \IncludeGraphicCentred{gabarito_aula4-1c}[0.9]
        \item \IncludeGraphicCentred{gabarito_aula4-1d}[0.9]
    \end{alphalist}
  \end{answer}
\end{exercise}

%------------------------------------------------------------------------------%
\begin{exercise}
  Encontre o domínio das seguintes funções:
  \begin{mathlist}
    \item g(x) = \frac{-5}{ \sqrt{16-\left( \frac{1}{2} \right)^x}}
    \item f(x) = \sqrt{9 \cdot 3^{2x-1}-\frac{1}{27}}.
  \end{mathlist}

  \begin{answer}
    \begin{mathlist}
      \item \{ x \in \R \st x    > -4 \}
      \item \{ x \in \R \st x \geq -2 \}
    \end{mathlist}
  \end{answer}
\end{exercise}

%------------------------------------------------------------------------------%
\begin{exercise}
  Resolva as equações abaixo
  \begin{mathlist}
    \item 8^{x} = 0,25
    \item 5^{2x-7} = 125
    \item 8^{x^2-x} = 4^{x+1}
    \item 2^{x^2- 3x - 4} = 1
    \item \sqrt{\left(\frac{1}{2}\right)^{3x-2}}
           = \left(\frac{1}{2}\right)^{-4x}.2^{-x+4}
    \item (10^x)^{1-x}-0,000001 = 0
    \item 2^{2x + 1} - 3.2^{x + 2} = 32
    \item 2^{x^2+5}+2^4 = 144
    \item 5^{4x-1}-5^{4x}-5^{4x+1}+5^{4x+2} = 480
    \item 2^{3x-2}-4^{x+6} = 0
  \end{mathlist}

  \begin{answer}
    \begin{mathlist}
      \item S = \left\{ -\frac{2}{3}        \right\}
      \item S = \left\{ 5                   \right\}
      \item S = \left\{ -\frac{1}{3}, 2     \right\}
      \item S = \left\{ -1, 4               \right\}
      \item S = \left\{ -\frac{2}{3}        \right\}
      \item S = \left\{ -2, 3               \right\}
      \item S = \left\{ 3                   \right\}
      \item S = \left\{ -\sqrt{2}, \sqrt{2} \right\}
      \item S = \left\{ \frac{1}{2}         \right\}
      \item S = \left\{ 14                  \right\}
    \end{mathlist}
  \end{answer}
\end{exercise}

%------------------------------------------------------------------------------%
\begin{exercise}
  Resolva as inequações abaixo
  \begin{mathlist}
    \item \left(\frac{1}{2}\right)^{2x-1}\leq \frac{1}{16}
    \item e^{x^2-5}\geq e^{4x}
    \item \left(\frac{1}{2} \right)^{2x-1}> \frac{1}{16}
    \item 7^{x-2}<343
    \item 1 \leq \left( \frac{1}{2}\right)^x \leq 8
  \end{mathlist}

  \begin{answer}
    \begin{mathlist}
      \item \left\{ x \in \R \st x \geq \frac{5}{2}            \right\}
      \item \left\{ x \in \R \st x \leq -1 \text{ ou } x\geq 5 \right\}
      \item \left\{ x \in \R \st x < \frac{5}{2}               \right\}
      \item \left\{ x \in \R \st x < 5                         \right\}
      \item \left\{ x \in \R \st -3 \leq x \leq 0              \right\}
    \end{mathlist}
  \end{answer}
\end{exercise}

%------------------------------------------------------------------------------%
\begin{exercise}
  Se \(f(x) = 16^{1+\sfrac{1}{x}}\),
  encontre o valor de $f(-1) + f(-2) + f(-4)$

  \begin{answer}
    13
  \end{answer}
\end{exercise}

%------------------------------------------------------------------------------%
\begin{exercise}
  Calcule $f(0)-f (3/2)$ sabendo que
  \[
    f(x) = \begin{cases}
      2^x,         & \text{ se } -1 \leq x \leq 1, \\
      \frac{1}{x}, & \text{ se } x > 1.
    \end{cases}
  \]

  \begin{answer}
    \(\frac{1}{3}\)
  \end{answer}
\end{exercise}

%------------------------------------------------------------------------------%
\begin{exercise}
  Seja a função $f\colon \R\to\R$ definida por $f(x) = 2^x$.
  Encontre o valor de $f(a+1)-f(a)$

  \begin{answer}
    $2^a$
  \end{answer}
\end{exercise}

\end{exercises*}

%------------------------------------------------------------------------------%

%------------------------------------------------------------------------------%
